% year: תש"פ
% type: מבחן
% semester: א
% moed: ג
% part: ב
% number: 4
% points: 9
% categories: func-analysis, derivatives, multiple-choice
% source: exams/מבחנים/תשפ/מועד מיוחד תשפ.pdf

\begin{question}
מצאו את הערך הקטן ביותר בתחום הגדרה של הפונקציה: $f(x)=x^4+4x$
\begin{enumerate}
  \item $-4$
  \item $1$
  \item $-3$
  \item $-5$
  \item $-2$
\end{enumerate}
\end{question}

\begin{solution}
\textbf{התשובה הנכונה: ג ($-3$).}

$f$ פולינום, מוגדר וגזיר בכל $\R$. $f'(x)=4x^3+4=4(x+1)(x^2-x+1)$, והגורם $x^2-x+1$ חיובי תמיד (דיסקרימיננטה שלילית). לכן $f'<0$ עבור $x<-1$ ו-$f'>0$ עבור $x>-1$: $f$ יורדת ב-$(-\infty,-1]$ ועולה ב-$[-1,\infty)$. מכאן ש-$x=-1$ היא נקודת מינימום מוחלט, והערך הקטן ביותר הוא
\[ f(-1)=1-4=-3 . \]
\end{solution}
